\section{Introduction}
Fine-grained visual recognition---telling apart 200 visually similar bird
species---is a setting where \emph{why} a model chose a particular class can
matter as much as the choice itself. Concept Bottleneck Models
(CBMs)~\cite{koh2020concept} tackle this by routing the prediction through an
intermediate layer of human-readable concepts: an image is first mapped to a
vector of concepts (the class-level bird attributes of CUB-200-2011%
~\cite{wah2011caltech}, selected following \citet{koh2020concept} and directly
supervised), and only then to a species label. This design exposes concept-level
explanations and allows test-time intervention on individual concepts.

Two problems, however, remain open in the standard CBM. First, concepts are
predicted from the \emph{whole} image. Because concepts and classes are
correlated in the training data, the encoder can pick up a concept such as
\emph{wing color} from features that are diagnostic of the \emph{species} rather
than of the wing itself. The concept vector then leaks class information and the
explanation becomes \emph{circular}: the concept is ``predicted'' from a
representation that already encodes the answer, so its activation no longer
reflects the local evidence it is supposed to describe, and concept interventions
lose their intended effect. Second, the label space is flat and the model is
evaluated by accuracy alone, treating every misclassification as equally
severe---confusing two warblers is penalised exactly like confusing a warbler
with an albatross.

We address both problems in a single framework. Specifically, we (i)~predict
each concept from a crop of the body part it describes, localised by
\emph{Keypoint R-CNN}, so that the concept must be grounded in its own visual
evidence; (ii)~add a coarse-to-fine structure to color and pattern concepts;
(iii)~equip the label space with a WordNet-based severity metric and a post-hoc
decision rule; and (iv)~analyse the resulting accuracy, severity and
interpretability trade-offs on CUB-200-2011.
